\chapter{自然选择不是“努力就会进化”}

\begin{kaichang}
打开家里的药箱，很可能躺着一板抗生素：阿莫西林、头孢，或者别的什么。翻翻说明书，几乎一定能找到这样一行字：“即使症状好转，也请完成整个疗程。”医生也会反复叮嘱：退烧了、嗓子不疼了，药也不能自己停。

这很奇怪。病都好了，为什么还要继续吃药？有人猜：是药厂想多卖药。有人猜：是怕病根没除干净，会复发。这两个猜测听起来都有点道理，但都不是真正的答案。真正的答案藏在一场已经进行了几十亿年、此刻可能正在某个病人体内上演的“筛选”里。先记下你的猜测，这一章结束时我们对答案。

还有一个更大的问题等着你。很多人这样讲进化：长颈鹿的脖子，是祖先们“为了够到高处的树叶，代代使劲伸，越伸越长”。也就是说，生物只要努力，就会朝着需要的方向进化。这个念头几乎每个人都动过——而它几乎每一部分都是错的。错在哪里？我们从那板抗生素说起。
\end{kaichang}

\section{一种正在失效的发明}

先看一组看得见、测得到的事实。

1928 年，弗莱明注意到培养皿里青霉菌周围的细菌死光了，由此发现了青霉素。到 1940 年代，青霉素实现量产，曾经在战场上致命的伤口感染忽然变得可治，肺炎、产褥热的死亡率应声而落。抗生素是二十世纪最伟大的医学发明之一，全球人均寿命的上升里有它实打实的一份功劳。

可是几十年后，事情起了变化。1960 年代，出现了对甲氧西林耐药的金黄色葡萄球菌（MRSA），这种“超级细菌”能在医院的床单、门把手上存活，普通抗生素拿它没办法。结核病的治疗也变了样：普通结核吃六个月药大多能治好，而对多种药物耐药的结核，疗程要拖到一年半甚至更久，药物副作用更重，治愈率还更低。世界卫生组织估计，2019 年全球有一百多万人的死亡与细菌耐药直接相关——注意这是模型估计，不是逐人清点的数字，但量级足够吓人。

关键的问题来了：这几十年里，抗生素的化学分子一个原子都没变过。青霉素的分子还是那个分子，杀菌的原理还是那个原理。\textbf{药没有变差，是细菌变了。}细菌不会开会商量，不会发愤图强，它们是怎么“变”的？

\section{细菌身上早就存在着差异}

把镜头推进到粒子层。你在第 72 章已经见过：细菌的遗传物质是一条环状 DNA，细胞分裂前要把整条链复制一遍。复制很精确，但不是完美——碱基配对偶尔会出错，修复系统能抓回绝大多数错误，可总有漏网的。这就是突变。

先感受一下数量级。条件合适时，一个大肠杆菌大约 20 分钟就能复制一次。你的肠道里常年住着大约 $10^{13}$ 个细菌。每个细胞每复制一次，平均大约会产生 $10^{-3}$ 个新突变。单个数字看着小，乘上天文数字般的个体数和繁殖速度，结论就反过来了：\textbf{在任何一个庞大的细菌群体里，几乎所有能发生的单点突变，此刻都正发生在某个细菌身上。}

这些突变里，绝大多数是中性的——改了一个不影响蛋白质工作的碱基，或者改在非关键位置，细菌照常过日子；一部分是有害的，让细菌长得慢、甚至活不成；极少数，恰好能挡住某种抗生素：有的突变改变了抗生素要攻击的那个蛋白质的形状，让药物分子对不上号；有的让细胞壁上多出一台“水泵”，把进来的药物分子再泵出去。这样的细菌，在千万个同伴里可能只有一两个。

现在轮到抗生素登场。药一下去，$99.9\%$ 的细菌死掉——包括所有没有那个突变的个体。而那一两个恰好带着耐药突变的细菌活了下来，面前是空出来的地盘和养分。它们复制、分裂，二十分钟后两个变四个，四个变八个……几天之内，一个全新的细菌群体建立起来，这一次，\textbf{每一个成员都是幸存者的后代，每一个都天生耐药}。

看清楚这台机器的逻辑：突变不是细菌“为了”抗药而制造的。突变的产生与药物毫无关系——早在抗生素出现之前几十亿年，突变就一直在发生；早在人类使用青霉素之前，土壤里就存在带耐药基因的细菌。抗生素做的只有一件事：\textbf{杀死不适者，留下碰巧合适者}。它不提供方向，它只是筛子。

\begin{zhengju}
突变究竟发生在药物出现之前，还是被药物“诱导”出来的？1943 年，卢里亚和德尔布吕克用一个漂亮的实验回答了这个问题。他们把同一批大肠杆菌分成许多小管培养，再把每管细菌分别涂到铺满噬菌体（一种专门感染细菌的病毒）的培养皿上，看各管能长出多少“抗噬菌体”的幸存者菌落。

如果抗性是被噬菌体“诱导”的，那么每个细菌遭遇噬菌体时获得抗性的概率相同，各管长出的菌落数应该差不多，像掷骰子的均匀分布。如果抗性突变是随机地、或早或晚发生在接触噬菌体之前，结果就完全不同：某管如果“运气好”，在早期就出了一个抗性突变，这个突变细菌繁殖了一代又一代，到涂皿时已经留下成千上万个后代——这管就会爆出几百个菌落；而多数管没赶上早的突变，菌落寥寥。结果正是后者：各管菌落数差异巨大，呈“头奖式”的分布。突变早于筛选，随机先于需要。这个思想实验与实验设计的结合，后来为他们赢得了诺贝尔奖，也为“突变不定向、选择定向”这条进化铁律钉上了第一颗钉子。
\end{zhengju}

\section{三个条件，一台机器}

现在可以把这台机器拆成零件了。自然选择要能运转，只需要三个条件同时成立：

\begin{itemize}[itemsep=2pt]
\item \textbf{可遗传的差异}：群体里的个体彼此不同，而且这些差异能通过 DNA 传给后代（第 72、73 章讲的就是差异从哪里来）。
\item \textbf{生存与繁殖的差异}：在\emph{当前}环境里，某些差异让携带者活得更久、留下更多后代。注意是“繁殖成功”，不只是“活着”——活得久却没后代，在进化账本上等于零。
\item \textbf{足够的时间}：一代接一代，有利差异的携带者在群体中占比越来越高。
\end{itemize}

三个条件都是纯机械的条件，没有任何一条提到“愿望”“努力”或者“进步”。差异是随机产生的，筛选是环境自动执行的，剩下的只是算术：谁的副本多，下一代里谁的声音就大。

为了说清楚“声音大”，我们需要一个符号。把一个基因的不同版本叫做\textbf{等位基因}（第 71 章）：耐药是一个版本，不耐药是另一个版本。某个等位基因在群体全部副本中所占的比例，叫做它的\textbf{等位基因频率}。进化的定义可以精确到一句话：\textbf{群体中等位基因频率随代际发生的改变}。频率变了，进化就发生了；跟“进步”“高级”一概无关。

\begin{qiaoliang}
这台机器有多快？算一笔账。假设一次感染里有一百万个细菌，其中恰好 1 个带着耐药突变——频率是 $10^{-6}$，百万分之一。抗生素杀死 $99.9\%$ 的敏感细菌：999999 个敏感者只剩大约 1000 个。那个耐药细菌活了下来。

接下来，幸存者在无人的地盘上繁殖。假设一个疗程过后、免疫系统收尾之前，每个幸存细菌平均留下 1000 个后代：敏感幸存者变成约 $10^6$ 个，耐药者变成约 $10^3$ 个。等等——这看起来耐药菌还是少数？别忘了第二轮：如果此时再来一轮抗生素（或者感染转到下一个人身上再遇药），敏感的这 $10^6$ 个又被砍到 $10^3$，而耐药的 $10^3$ 个毫发无损、继续扩增到 $10^6$。两轮之后，耐药等位基因的频率从百万分之一变成约 $99.9\%$。

从几乎为零到几乎垄断，只用了两轮筛选。这就是医生最怕的算术，也是“进化可以很快”的直接证据：不需要几百万年，只要差异够要命、繁殖够快，几天就够了。
\end{qiaoliang}

\section{蛾子、鸟喙和风向标}

细菌藏在培养皿里，看不见摸不着。好在自然选择的三个条件在任何生物身上都成立，包括看得见的动物。两个经典案例，值得一个个讲。

\subsection{桦尺蛾：写在树皮上的频率}

英国有一种桦尺蛾，十九世纪上半叶的采集记录里几乎全是浅色型——翅膀灰白带黑斑，停在覆满地衣的浅色树皮上，鸟很难发现。1848 年前后，曼彻斯特附近第一次记录到黑色型个体；到 1895 年，当地黑色型的比例已经超过 $90\%$。五十年间，一个等位基因的频率从稀有变成主流。为什么偏偏是在曼彻斯特——工业革命烟囱最密集的城市？

解释链条是这样的：煤烟熏死了树皮上的浅色地衣，又把树干熏黑。浅色蛾停在黑树干上，从“隐身”变成“显眼”，被鸟啄食的概率上升；黑色型反过来占了便宜。注意因果的方向：煤烟没有“把蛾子熏黑”，单只蛾子的颜色一辈子不会变；变的是\emph{群体里两种颜色各自留下后代的机会}。

\begin{zhengju}
这个解释真的对吗？会不会是煤烟直接改变了蛾子的基因？1950 年代，生态学家凯特威尔设计了实验：他在污染的乡村和未污染的林地分别释放做了标记的两种蛾子，再设法回收。结果：在污染区，黑色型的回收率显著高于浅色型；在未污染区，正好反过来。他还直接观察到鸟确实叼走显眼的个体。

后来有批评者指出凯特威尔实验的漏洞：桦尺蛾在自然状态下大多停在树冠高处，而不是他放置蛾子的树干上，实验装置可能夸大了差异。这是科学该有的样子——实验设计被同行挑刺，结论暂时打上问号。2000 年代，马杰勒斯用更贴近自然的方式重做：把蛾子释放在树上它们真正停留的位置，多年观察。结果再次确认：鸟的捕食确实按体色区别对待，差异足以解释频率变化。与此同时，英国治理空气污染之后，树皮重新变亮，黑色型的频率果然回落——\textbf{预言兑现}：环境翻回去，选择方向就翻回去。一个理论能预言未来的变化方向，才算真的抓住了因果。
\end{zhengju}

\subsection{达尔文雀：在眼皮底下的逆转}

如果说蛾子的故事发生在五十年尺度上，加拉帕戈斯群岛的地雀把进化压缩到了几年。1977 年，格兰特夫妇正在岛上逐年测量中地雀的喙。这一年大旱，小而易剥的种子先被吃光，只剩又大又硬的种子。喙小而弱的鸟嗑不开硬种子，大量饿死；喙稍深、稍有力的鸟活下来并繁殖。第二年孵出的幼鸟，平均喙深比干旱前大了约 $4\%$——只是几十分之一的差别，但方向清晰可测。

故事没有停在“喙越来越大”上。1982 到 1983 年，厄尔尼诺带来暴雨，小草疯长，小种子重新泛滥。这一回轮到喙大而笨的鸟吃亏——对付小种子它们效率低下。几季之后，平均喙深又降了回去。\textbf{同一个群体，几年之内被选择先推向一边、再推回另一边。}没有永恒的“更好”，只有“此时此刻更合算”。达尔文当年只看到群岛之间鸟喙的静态差异，而他的后继者在同一座岛上看到了活的、会反转的进化。

\subsection{回到药瓶}

抗生素耐药是同一台机器的第三个例子，也是最贴近你的一个。而且它和雀喙一样，藏着反转的线索：许多耐药突变是有代价的。那台把药物泵出去的“水泵”要消耗能量，改变形状的蛋白质可能干活变慢——在没有抗生素的环境里，耐药细菌常常竞争不过利索的敏感菌，频率会慢慢回落。这就是为什么滥用抗生素（感冒也吃、吃两天就停、给牲畜当促生长剂喂）是全人类的损失：它让“筛子”时时刻刻悬在细菌头上，把耐药型一代代保留下来。

\section{谁在变，谁不变}

现在可以拆除那个流传最广的误会了。把三个案例并排看，你会发现同一件事：

\begin{itemize}[itemsep=2pt]
\item 没有哪一只蛾子“变成了”黑色；是黑色型在群体中的\emph{频率}变了。
\item 没有哪只地雀的喙在干旱里“长长”了；是喙深的鸟留下的后代多，下一代的平均值移动了。
\item 没有哪个细菌“练出”了抗药性；是碰巧耐药的个体活了下来，独占了下一代。
\end{itemize}

单一个体的一生，走的是\textbf{发育}：一个受精卵按基因调控网络的程序长大、成熟、衰老（第 75 章），基因型基本不变。进化不发生在个体身上，它发生在\textbf{群体}的代与代之间，变的是等位基因的频率。\textbf{个体发育，群体进化。}把这两件事混为一谈，是“努力就会进化”这种错觉的总根源。

长颈鹿的故事也该重讲了。不是祖先们“使劲伸脖子，把伸长的脖子传给孩子”。真实的过程是：鹿群里本来就存在脖长的可遗传差异；在低矮食物稀缺的季节，脖子略长的个体吃到更多叶子、留下更多后代；许多代之后，整个群体的平均脖长移动了。没有任何一只鹿努力过，也没有任何一只鹿需要努力——筛选自动完成了一切。

\begin{wuqu}
“用进废退，获得性遗传。”这是拉马克在达尔文之前提出的著名猜想：铁匠胳膊粗，儿子天生胳膊粗；长颈鹿伸脖子，后代脖子就变长。它错在把\emph{个体的变化}当成了\emph{可遗传的变化}。肌肉练得再发达，生殖细胞里的 DNA 序列也不会照着肌肉的样子改写——举重运动员的孩子出生时和别的孩子一样，得自己练。反例随手可得：几代人都戴耳环的家庭，女婴出生时耳垂上并没有孔。

顺带说一句公道话：第 74 章讲过，某些表观遗传状态确实能跨代传递，环境也真能影响发育中的性状（第 73 章）。但这些现象既不“定向”——生物无法按需定制改变——也没有推翻主线：可遗传变异的根本来源仍是 DNA 层面的变化，频率的改变仍靠群体层面的筛选。拉马克猜错了机制，达尔文（以及后来的遗传学）给出了对的那一个。
\end{wuqu}

\section{适应只对“当时当地”负责}

三个案例还教会我们另一件事：\textbf{适应没有方向指向未来，它只对当前环境结账}。黑色型蛾子在污染年代是“优等生”，空气一干净立刻变回“差等生”；大喙在旱年是救命的，涝年就成了累赘；耐药细菌在有药的世界里称王，药一撤就把地盘还给敏感菌。进化里没有永远的赢家，因为“赢”的定义写在环境里，而环境一直在变。

所以“适应”不等于“完美”，只等于“此刻够用，而且比对手强一点”。自然选择是个得过且过的工程师：它不从零设计，只在现有变异里挑挑拣拣，够用就过关。脊椎动物的眼睛里，视神经穿过视网膜的地方留下一个没有感光细胞的盲点；人的气管和食管在咽喉处交叉，使“噎住”成为一种可能的事故。这些“设计缺陷”不是进化的失败记录，恰恰是它的签名——只有“在旧图纸上修修补补、够用就好”的过程，才会留下这种痕迹。一个从零设计的工程师不会这样干活。

这也意味着进化没有终点、没有阶梯。不是从“低等”走向“高等”，更不是朝着人类攀登（第 78 章会展开这棵树的真实形状）。细菌被筛选了几十亿年，今天依然是细菌——而且在它们的世界里适应得极好。人类不是进化的目的，只是无数支系中的一支。

\begin{shiyan}
\textbf{模拟：当一回“鸟”，亲手筛一次}

材料：一张花色复杂的包装纸（或花桌布）铺桌面；两种颜色的纸屑各 50 片——一种颜色接近桌布底色，一种反差强烈；镊子一把；计时器。

步骤：①把两种纸屑混匀撒在桌布上。②请一位同学当“捕食者”，用镊子逐片夹取，限时 20 秒，夹到一片就放一边。③数一数两种颜色各被夹走多少片。④让“幸存者”繁殖：每剩下一片，就补一片同色的进去，凑回 100 片。⑤重复②到④，共三到四轮，记录每轮两种颜色的比例变化。

观察点：哪一种颜色的比例逐轮上升？如果你的“捕食”足够随机而环境（桌布）不变，频率会稳定地朝一个方向移动——你刚刚亲手执行了一次自然选择。换一张底色不同的桌布再玩一遍，方向会改变吗？

安全提示：使用儿童剪刀裁纸屑，镊子不要对着人；纸屑游戏请勿入口。
\end{shiyan}

\section{回到那板抗生素}

现在兑现开场的承诺：为什么症状消失了还必须把疗程吃完？

因为细菌群体对药物的敏感性不是整齐划一的，而是一条连续谱：最敏感的个体先死，抵抗力稍强的能撑更久。你吃下头几天药，体温退了、嗓子不疼了——那是因为细菌的主力部队已被消灭，数量低到免疫系统能轻松应付。但此时活下来的，恰恰是群体里\emph{最耐药的那一小撮}。如果此时停药，等于在筛选进行到一半时撤走了筛子：这些“预选赛优胜者”获得喘息之机，重新繁殖，建起一个平均耐药水平更高的新群体。你不仅可能复发，还可能把升级过的细菌传染给下一个人。

完整的疗程，是用足够长的药物暴露时间把筛选进行到底，配合免疫系统把幸存者清到零。\textbf{“提前停药”不是在省药，是在亲手训练更耐药的细菌。}当然，具体到每一次生病该不该用抗生素、用哪种、用多久，要交给医生判断——病毒性感冒吃抗生素毫无用处，只会白白筛选你肠道里的常驻菌群。这属于医嘱范畴，本书只负责讲清背后的进化算术。

同一套算术也在农田里运行。一种新杀虫剂投入使用的头几年往往效果显著，几年后却越来越不灵——不是药剂“失效了”，而是害虫群体被筛选出了抗性。DDT 的历史就是完整的一轮：1940 年代它几乎是灭蚊神药，十年之内，多个蚊种里抗性个体的频率从稀有变成主流。农业上如今的对策也带着鲜明的进化眼光：轮换使用机制不同的药剂，让“筛子”的方向不断变化；或者在田边保留一片不喷药的“庇护所”，让敏感个体活下来、持续稀释抗性等位基因的频率。理解了“选择只是筛选”，这些策略就不再是口诀，而是必然推论。

\section{这台机器何时不灵，以及之后的故事}

自然选择很强大，但请记住它的边界，知道它什么时候“转不动”：

\begin{itemize}[itemsep=2pt]
\item 群体里如果没有相关的可遗传差异，再强的筛选也无米下锅——这就是为什么多样性本身就是群体的保命资本。
\item 如果差异不影响生存和繁殖，选择对它们“视而不见”，这些等位基因的频率照样会变，但靠的是纯粹的运气——遗传漂变。
\item 如果群体很小，运气甚至可以压倒实力：再好的等位基因，也可能因为携带者“没赶上繁殖”而消失。
\end{itemize}

后两种情况不是本章的例外补丁，而是另一台机器的存在证明：进化不只有自然选择。突变供给原料、漂变掷骰子、基因流在群体间搬运等位基因——第 77 章就来讲这些“非选择”的力量，以及它们怎样让分子层面的变化变成一只可以读数的“时钟”。而当一个群体被山川海洋隔开、各自积累变化，最终会怎样分成两个物种？那是第 78 章的故事。

\begin{tiaozhan}
（跳过不影响主线）选择推动频率变化，能不能写出公式？可以。设某等位基因的频率为 $p$，携带它的个体平均每代留下 $w_1$ 个后代（这叫\textbf{适合度}），携带另一版本的个体适合度为 $w_2$。令相对适合度 $w_1 = 1$，$w_2 = 1 - s$，其中 $s$ 叫\textbf{选择系数}：$s = 0.1$ 意味着该版本每代平均少留下 $10\%$ 的后代。经过一代，频率近似变为
\[
p' = \frac{p}{p + (1-p)(1-s)}.
\]
代入耐药菌的例子：$p = 10^{-6}$，有药环境下敏感型的 $s$ 接近 $0.999$，你算算一代之后 $p'$ 是多少？再代入桦尺蛾：估算显示黑色型的优势相当于 $s$ 约为 $0.3$ 量级——看起来温和的优势，五十代（五十年，桦尺蛾一年一代）足以把频率从 $1\%$ 推到 $90\%$ 以上。你可以用计算器验证：进化不快不慢，它正好是复利。
\end{tiaozhan}

\section*{练一练}

\begin{sikaoti}
\item 用“个体发育、群体进化”重新表述这句话，指出原句错在哪：“这只地雀在干旱中努力把喙长得更大，以适应坚硬的种子。”
\item 画一幅信息流图：从“DNA 复制出错”到“群体耐药频率上升”，标出哪一步是随机的、哪一步是有方向的，并说明方向由谁提供。
\item 在一个没有抗生素的环境里，耐药细菌的频率往往会缓慢下降。用“适应只对当时环境负责”和“耐药的代价”解释这个现象，并预测：如果此时重新用药，频率会怎样变化？
\item 有报道说“细菌为了抵抗抗生素而进化出耐药性”。这句话至少有三处用词会误导人，把它们找出来并改写成准确的说法。
\item 假设某岛屿引入了一种新的捕食者，只捕食飞行中的昆虫。请预测：若干代之后，这个昆虫群体可能在哪些性状上发生频率变化？你的预测依赖哪三个条件成立？
\item 参考 qiaoliang 中的复利算术：某种等位基因每代频率翻倍的情况不存在，但“有利者留下更多后代”与银行存款复利有何相似之处？这个类比在哪里会失效？
\end{sikaoti}
